\section{Introduction}

Intra-urban population surfaces are useful because many planning and analytic questions are spatial before they are causal. Service screening, hotspot-aware prioritization, exploratory comparison across neighborhoods, and first-pass city diagnostics often need a gridded view of where population is more likely to concentrate rather than only a single total for the whole city. Large-scale gridded population products and dasymetric approaches have shown the value of redistributing population counts with spatial covariates, especially where fine-grained census geography is limited or unevenly available \cite{Leyk2019,Mennis2009,Stevens2015}.

At the same time, deterministic surfaces can be misleading in data-scarce settings. A map that assigns one number per cell can look equally trustworthy everywhere even when local evidence quality differs sharply across cities and within cities. This is a practical problem for proxy population work: when cell-level census labels are unavailable, analysts may still need a bounded screening layer, but they also need to know where the resulting surface is more stable and where it should be treated more cautiously.

City1 v2 addressed the first half of this problem by building a calibrated proxy population surface for Kazakhstan from a frozen 500 m workflow with official-total calibration, OpenStreetMap-derived predictors, weak supervision, and multi-level validation. Its identity was explicit: proxy, not truth. That deterministic baseline remains useful, but it leaves open a second question that matters for reviewer scrutiny and downstream interpretation: how should users distinguish stronger and weaker parts of the calibrated proxy surface when cell-level truth is unavailable?

City1 v3 is designed as a bounded extension rather than a scientific reset. It keeps the v2 calibrated proxy-surface identity and adds an uncertainty-aware interpretation layer around it. The goal is not to estimate true census uncertainty, and not to replace the v2 surface with a new task identity. The goal is to quantify how stable the proxy surface appears under a frozen ensemble configuration and to expose that stability through intervals, confidence bands, and hotspot classes that can support more careful screening.

The contributions of City1 v3 are:

\begin{itemize}
\item uncertainty intervals for calibrated proxy population surfaces through ensemble-based P10/P50/P90 outputs;
\item confidence bands for reliability-aware interpretation, with an explicit statement that \texttt{confidence\_score} is not a truth probability;
\item hotspot prioritization that separates stable screening candidates from caution-heavy ones;
\item a bounded uncertainty-validation stack covering interval behavior, error-width alignment, confidence-band separation, hotspot stability, district aggregation limits, and external structural comparison;
\item a frozen paper-facing evidence package that makes all Phase 8 manuscript claims traceable to fixed tables, figures, manifests, and limitation files.
\end{itemize}

Under this framing, City1 v3 should be read as an uncertainty-aware screening and interpretation layer for data-scarce urban settings. Its strongest value is not that it proves true population uncertainty, but that it reduces deterministic over-reading of a calibrated proxy surface while preserving the transparency of the v2 baseline.

